\section{总结}
本作品针对传统加密技术难以隐藏通信行为、部分深度学习图像隐写方案需要额外训练专用网络等问题，设计并实现了基于图像密钥控制DNN状态的安全隐蔽信息表示框架SSIR。
该框架直接利用现有DNN的内部神经元状态表示消息，通过图像密钥和对抗优化实现信息的生成与恢复，无需修改模型参数或原有功能。主要完成了以下四个方面的工作：

1. 提出了基于DNN状态的信息表示方法。 将消息编码为选定神经元的目标状态，通过图像密钥融合与对抗优化生成密文图像，接收方利用相同密钥和DNN恢复原始消息。

2. 设计了分组编码机制。 针对部分神经元对输入扰动不敏感的问题，通过对相邻神经元响应取平均，减少单个神经元不响应的影响，提高目标状态诱导的可靠性。

3. 设计了双重信息泄漏抑制方法。 针对密文图像可能存在的直接视觉泄漏和间接DNN状态泄漏，设计联合损失函数，在生成过程中同时减少视觉变化和神经元状态偏差。

4. 进行了形式化安全分析。 分析了密文图像的概率生成与多语义特性，推导了查询有界的IND-CPA和IND-CCA1区分优势上界，并给出了查询有界安全性所需的分布条件和单密钥查询次数限制。

本作品在ImageNet数据集和真实图像上进行了实验测试。结果表明，分组编码将所测生成成功率由79\%提高至100\%；通过调整优化参数，直接视觉泄漏指标最高降低99.79\%，
启用间接泄漏约束后，神经元状态偏差平均降低56.02\%。对于150位消息，密文图像生成时间约为300～500 ms，消息恢复时间约为10 ms。SSIR还能够适配多种常见DNN，并在测试条件下表现出一定的攻击鲁棒性和图像失真容忍能力。

综上，SSIR在不额外训练专用网络的情况下，实现了由图像密钥控制的信息表示与恢复，并通过分组编码和泄漏抑制机制提高了方案的
可靠性与隐蔽性。该方案可用于进一步探索会话密钥传递、物联网控制指令和企业内部短消息等信息的隐蔽传输，但实际部署时仍需考虑密钥共享、查询次数限制和密钥更新等问题。